\section{实验环境}\label{sec:task2-environment}

软件环境取自本次工作区可直接检查到的运行环境。Python 解释器使用
\path{/Users/jiaqiaosu/miniconda3/bin/python}；该解释器能够导入
\texttt{cv2}，版本由命令 \texttt{cv2.\_\_version\_\_} 确认。硬件和采集参数不根据软件环境推断，须在实际实验后补录。

\begin{table}[htbp]
    \centering
    \caption{任务 2 实验环境与待补录采集信息}
    \label{tab:task2-environment}
    \begin{tabular}{lll}
        \hline
        类别 & 项目 & 记录值 \\
        \hline
        操作系统 & 系统与架构 & macOS 26.6.2，Darwin 25.6.0，arm64 \\
        软件 & Python & 3.13.12 \\
        软件 & OpenCV & 4.10.0（\texttt{cv2}） \\
        软件 & 主要接口 & \texttt{findChessboardCorners}、\texttt{cornerSubPix} \\
        软件 & 主要接口 & \texttt{calibrateCamera}、\texttt{projectPoints} \\
        采集硬件 & 手机型号 & 待实际采集/待填写 \\
        采集硬件 & 使用摄像头 & 待实际采集/待填写（前置或后置） \\
        图像 & 图像分辨率 & 待实际采集/待填写 \\
        标定板 & 内角点规格 & 待实际采集/待填写 \\
        标定板 & 方格实际边长 & 待实际采集/待填写 \\
        数据 & 采集图像总数 & 待实际采集/待填写 \\
        \hline
    \end{tabular}
\end{table}

表\ref{tab:task2-environment}中的软件版本已经由本机命令核实；其余字段是实验数据字段，不能用默认配置或估计值替代。正式提交前，应将手机系统相机设置（如有改变分辨率、裁剪或数码变焦）一并记录。
